\documentclass[11pt]{article}

\usepackage[utf8]{inputenc}
\usepackage[T1]{fontenc}
\usepackage{amsmath, amssymb}
\usepackage{geometry}
\usepackage{hyperref}

\geometry{margin=1in}

\title{Photons and the Geometry of Spacetime:\\
Reasoning About Photon Reference Frames}
\newcommand{\authorname}{Reivax Del Ponte}
\newcommand{\institute}{Pseudo-Institute of Non-Experimental Physics}
\author{\authorname \thanks{\institute}}

\date{\today}

\begin{document}
\maketitle

\begin{abstract}
We examine the consequences of taking the reference frame of a single photon seriously, treating emission and absorption as coincident events in that degenerate limit, and explore how this perspective informs our reasoning about multiple photons, the structure of spacetime, and the topology of causal sequences. We argue that, while the photon's proper frame collapses time and distance, the topology of spacetime is preserved, and a consistent picture of how photons and event sequences are to be reasoned about within that structure can be reconstructed.
\end{abstract}

\section{Introduction}

Special and general relativity provide a rich framework for reasoning about light, photons, and the geometry of spacetime. One question that is rarely entertained in standard treatments, however, is \emph{what does the world look like in the reference frame of an individual photon?} The reason is technical: a photon has no rest mass, and the standard passage to a rest frame via a Lorentz boost produces a singular limit. We argue that this limit, taken seriously rather than discarded, clarifies several subtle issues surrounding the topology of spacetime, the relationship between emission and absorption events, and the structures by which to reason about multiple photons.

\section{The Photon Reference Frame as a Degenerate Limit}

Consider a photon emitted at spacetime event $E = (t_e, \vec{x}_e)$ and absorbed at event $A = (t_a, \vec{x}_a)$. The photon's worldline connects these two events, and the spacetime interval between them is, by definition, lightlike:
\begin{equation}
s^2 \;=\; -c^2(t_a - t_e)^2 + \|\vec{x}_a - \vec{x}_e\|^2 \;=\; 0,
\label{eq:lightlike}
\end{equation}
so that $\|\vec{x}_a - \vec{x}_e\| = c \, (t_a - t_e)$, as is appropriate for a particle travelling at the speed of light.

\subsection{Lorentz Boost Along the Photon Worldline}

A Lorentz boost with velocity $v$ along the direction of photon propagation transforms the temporal and spatial separations as
\begin{align}
\Delta t' \;=\; \gamma\!\left(\Delta t - \frac{v\,\Delta x}{c^2}\right),
\\[4pt]
\Delta x' \;=\; \gamma\,(\Delta x - v\,\Delta t),
\end{align}
where $\gamma = 1/\sqrt{1 - v^2/c^2}$ is the Lorentz factor. Substituting the photon equation of motion $\Delta x = c\,\Delta t$ yields
\begin{align}
\Delta t' \;=\; \gamma\,\Delta t\!\left(1 - \frac{v}{c}\right),
\label{eq:deltatprime}
\\[4pt]
\Delta x' \;=\; \gamma\,\Delta t\,(c - v).
\label{eq:deltaxprime}
\end{align}

To examine the limit $v \to c$, set $v = c(1 - \varepsilon)$ for small $\varepsilon > 0$. Then
\[
\gamma \;\sim\; \frac{1}{\sqrt{2\varepsilon}},
\qquad
1 - \frac{v}{c} \;=\; \varepsilon,
\]
so that
\begin{align}
\Delta t' &\;\sim\; \frac{\Delta t \cdot \varepsilon}{\sqrt{2\varepsilon}}
\;=\; \Delta t\sqrt{\frac{\varepsilon}{2}} \;\longrightarrow\; 0,
\\[4pt]
\Delta x' &\;\sim\; \frac{\Delta t \cdot c\varepsilon}{\sqrt{2\varepsilon}}
\;=\; \Delta t\, c\sqrt{\frac{\varepsilon}{2}} \;\longrightarrow\; 0.
\end{align}
Both the temporal and spatial separations between emission and absorption vanish identically in the limit: in the photon's ``rest frame'', the two events coincide.

\subsection{Causal Order Under Lorentz Boosts}

A natural worry is whether some inertial frame could see absorption occur \emph{before} emission. Equation~\eqref{eq:deltatprime} settles the question: for any $v < c$ (the maximum boost for an inertial frame) we have $\gamma > 0$, $\Delta t > 0$, and $1 - v/c > 0$, and therefore $\Delta t' > 0$. The order of emission and absorption is the same in every inertial frame, and the photon frame in particular blends them to a single event. Causal order is absolute on the light cone.

\subsection{Topology Preservation}

Lorentz transformations are isometries of Minkowski space and therefore preserve all topological properties of spacetime. The lightlike character of the interval \eqref{eq:lightlike} is invariant, and although the limit $v \to c$ collapses the metric distance between $E$ and $A$ to zero, the two events remain \emph{logically} distinct: the topology of the relevant region of spacetime is unchanged. The endpoints $E$ and $A$ are not erased, but rather identified by the limit geometry.

\section{Before and After the Photon}

Standard physics declares the photon reference frame to be ``invalid'' on the grounds that the limit $v \to c$ is singular. We challenge this position, and instead interpret the photon frame as an \emph{instantaneous relation} between emission and absorption: in that frame, both events occur at the same time and at the same place.

\subsection{What Happened ``Before'' and ``After''?}

In the photon's frame, there is no before and no after, since time does not flow. In any other inertial frame, however, there are instants at which the photon ``does not yet exist'' (before emission) and instants at which the photon ``no longer exists'' (after absorption). With the notion of ``existence'' in quotation marks, because whether a given observer can unambiguously say that the photon exists depends on whether their worldline intersects the photon's --- a question whose answer is itself frame-relative in a precise sense.

\subsection{The Midpoint Question}

Consider photons travelling from the Sun to the Earth. In the Earth's frame, a given photon was emitted $\Delta t$ minutes ago and will be absorbed now. Does it make sense to say that, $\Delta t/2$ minutes ago, the photon was halfway between Sun and Earth?

The answer depends on the chosen frame:
\begin{itemize}
\item In the photon's frame, every point on the worldline is coincident with emission and absorption. The notion of ``halfway'' has no meaning, because all intermediate spacetime points are identified with the endpoints.
\item In the Earth's frame (or any other non-degenerate frame), the photon's worldline is parameterised by coordinate time, and the midpoint is a perfectly well-defined spacetime event. Physical interactions at that event (blocking, scattering, amplification) are local and frame-invariant in their consequences.
\end{itemize}

That the midpoint is identified with the endpoints in the photon frame does not contradict its physical reality in other frames. Identification is a feature of one particular degenerate limit, not a denial that the worldline has intermediate structure.

\subsection{A Thought Experiment With a Blocking Plate}

Replace the Sun with a single-photon laser emitting one photon per second, and place a thin plate halfway between source and destination. Program the plate to block each photon for one-minute intervals, indexed by a known sequence (e.g., the prime numbers). Photons take $\Delta t$ minutes to cross the gap, so the receiver should observe the corresponding block/unblock pattern offset by $\Delta t/2$ minutes relative to the laser.

The general-relativistic machinery confirms the expectation. The geodesic equation governing the photon's trajectory,
\begin{equation}
\frac{d^2 x^\mu}{d\lambda^2} + \Gamma^\mu_{\alpha\beta}\,
\frac{dx^\alpha}{d\lambda}\,\frac{dx^\beta}{d\lambda} \;=\; 0,
\label{eq:geodesic}
\end{equation}
determines the worldline for any choice of coordinates, in any background geometry specified by the Einstein field equations
\begin{equation}
G_{\mu\nu} + \Lambda\,g_{\mu\nu} \;=\; \frac{8\pi G}{c^4}\,T_{\mu\nu}.
\label{eq:efe}
\end{equation}
The plate blocks a photon if and only if the photon's null geodesic intersects the plate's worldtube. This condition is frame-invariant, so the observed sequence at the receiver is in fact the prime-indexed pattern, appropriately delayed: the midpoint is a physically meaningful event, even though in the photon frame it is identified with emission and absorption.

\section{The Two Photons Problem}

Let $A$ and $B$ be two objects. Each emits a photon towards the other at the same coordinate time, producing photon $p_{AB}$ travelling from $A$ to $B$ and photon $p_{BA}$ travelling from $B$ to $A$. In any inertial frame, the two photons travel towards each other and meet at a single spacetime event $M$, the midpoint of the $A$--$B$ separation along the line joining them. At $M$, the photons coincide and may interfere --- a phenomenon with well-established experimental confirmation.

\subsection{Two Distinct Degenerate Frames}

The reference frames of $p_{AB}$ and $p_{BA}$ are different --- they correspond to different null directions. In $p_{AB}$'s frame, the photon is emitted and absorbed at the same spacetime event; the same holds in $p_{BA}$'s frame for that photon. Neither photon has a global rest frame in which both photons are simultaneously identified with their endpoints.

\subsection{Can $p_{AB}$ Know About $p_{BA}$?}

Within the degenerate frame of $p_{AB}$, the only spacetime event in which $p_{BA}$ can be said to ``exist'' is the meeting event $M$. Before $M$, $p_{AB}$ has not yet encountered $p_{BA}$; after $M$, $p_{AB}$ has already passed $p_{BA}$. From $p_{AB}$'s point of view, $p_{BA}$ exists only as a pointwise coincidence at $M$, and at no other spacetime point.

The sharp boundary is the following: in a degenerate photon frame, knowledge of another photon is restricted to the set of intersection events between the two worldlines. Outside those intersections, the other photon is identified with no particular spacetime event, and so cannot be located.

\section{The Sequence Problem}

Let us now consider two objects $A$ and $B$ and a sequence of two photon exchanges: first, $A$ emits a photon $p_{AB}$ towards $B$; once $p_{AB}$ is absorbed, $B$ emits a photon $p_{BA}$ towards $A$. In the degenerate frames of either photon, time does not flow, and the two exchanges are simultaneous. In any inertial frame, however, the following sequence holds:
\begin{enumerate}
\item Neither $p_{AB}$ nor $p_{BA}$ exists.
\item $A$ emits $p_{AB}$.
\item $p_{AB}$ exists; $p_{BA}$ does not.
\item $B$ absorbs $p_{AB}$.
\item Neither $p_{AB}$ nor $p_{BA}$ exists.
\item $B$ emits $p_{BA}$.
\item $p_{AB}$ does not exist; $p_{BA}$ exists.
\item $A$ absorbs $p_{BA}$.
\item Neither $p_{AB}$ nor $p_{BA}$ exists.
\end{enumerate}

Four spacetime events anchor the sequence: $A_e$ (emission by $A$), $B_a$ (absorption by $B$), $B_e$ (emission by $B$), and $A_a$ (absorption by $A$). The intervals
\[
A_e \to B_a, \qquad B_e \to A_a
\]
are lightlike, since each connects the endpoints of a photon worldline. The intervals
\[
A_e \to B_e, \qquad B_a \to A_a, \qquad A_e \to A_a, \qquad B_a \to B_e
\]
are timelike within small enough regions of spacetime, since emission and absorption of consecutive photons occur at the same object separated by a finite interval.

The topology of the underlying causal graph --- with vertices $\{A_e, B_a, B_e, A_a\}$ and lightlike edges $\{A_eB_a, B_eA_a\}$ --- is invariant under Lorentz transformation. The degenerate photon frame does not destroy this graph; it only identifies certain vertices (the endpoints of each edge) in its particular limit geometry. The chain $A_e \to B_a \to B_e \to A_a$ remains a chain, even though in $p_{AB}$'s frame $A_e$ and $B_a$ coincide and in $p_{BA}$'s frame $B_e$ and $A_a$ coincide.

\section{Conclusion}

Allowing the reference frame of a photon as a degenerate limit provides a useful tool for reasoning about the structure of spacetime. In the limit, time does not flow and distances vanish, but the topology of spacetime is preserved: objects $A$ and $B$ retain their identities as endpoints of photon worldlines, even when those worldlines are collapsed.

Within this framework, other photons exist in the structure of spacetime in a photon's reference frame, but only at the spacetime events where their worldlines intersect the photon's own worldline. Sequences of emission and absorption events can be reasoned about as a graph whose edges are null geodesics and whose vertices are spacetime events. The frame collapse identifies certain vertex pairs (the endpoints of each photon's worldline) but does not destroy the graph.

The key insight is that the photon frame is not a fully fledged coordinate system but rather a bookkeeping device: it identifies which events are causally related by lightlike separation, without supplying coordinates for those events. Reasoning about the structure of spacetime without coordinates is unusual but tractable, and the picture sketched here is internally consistent: emission precedes absorption in every inertial frame, the topology of causal sequences is preserved, and other photons exist in the structure only at intersection events.

\end{document}
